\section{More subtractive methods}

\subsection{Sums with varying signs}

\begin{definition}
\label{def.cancel.all-even}
\lean{AlgebraicCombinatorics.SubtractiveMethods.IsAllEven}
\leanhelper
\leanok
Let $n\in\mathbb{N}$ and $d\in\mathbb{N}$. An
$n$-tuple $\left(  x_{1},x_{2},\ldots,x_{n}\right)  \in\left[  d\right]  ^{n}$
is said to be \emph{all-even} if each element of $\left[  d\right]  $ occurs
an even number of times in this $n$-tuple (i.e., if for each $k\in\left[
d\right]  $, the number of all $i\in\left[  n\right]  $ satisfying $x_{i}=k$
is even).
\end{definition}

In the Lean formalization, sign tuples $\left(e_1, e_2, \ldots, e_d\right) \in \{1, -1\}^d$
are represented as functions $e : \operatorname{Fin} d \to \mathbb{Z}/2\mathbb{Z}$,
where $0$ represents $+1$ and $1$ represents $-1$. The conversion is performed
by a sign conversion function. The product $e_{x_1} e_{x_2} \cdots e_{x_n}$
is represented by a sign product function, and the sum $e_1 + e_2 + \cdots + e_d$
(in $\mathbb{Z}$) by a sign sum function.

\begin{definition}
\label{def.cancel.multiplicity}
\lean{AlgebraicCombinatorics.SubtractiveMethods.multiplicity}
\leanhelper
\leanok
For an $n$-tuple $x : \operatorname{Fin} n \to \operatorname{Fin} d$ and
$k \in \operatorname{Fin} d$, the \emph{multiplicity} of $k$ in $x$ is the
number of all $i \in \operatorname{Fin} n$ satisfying $x(i) = k$.
\end{definition}

\begin{definition}[Sign conversion]
\label{def.cancel.toSign}
\lean{AlgebraicCombinatorics.SubtractiveMethods.toSign}
\leanhelper
\leanok
The function $\operatorname{toSign} : \mathbb{Z}/2\mathbb{Z} \to \mathbb{Z}$
converts elements of $\mathbb{Z}/2\mathbb{Z}$ to signs:
$\operatorname{toSign}(0) = 1$ and $\operatorname{toSign}(1) = -1$.
\end{definition}

\begin{definition}[Sign product]
\label{def.cancel.signProduct}
\lean{AlgebraicCombinatorics.SubtractiveMethods.signProduct}
\leanhelper
\leanok
For a sign tuple $e : \operatorname{Fin} d \to \mathbb{Z}/2\mathbb{Z}$ and an
index tuple $x : \operatorname{Fin} n \to \operatorname{Fin} d$, the
\emph{sign product} is
\[
\operatorname{signProduct}(e, x) := \prod_{i=1}^{n} \operatorname{toSign}(e_{x_i}).
\]
\end{definition}

\begin{definition}[Sign sum]
\label{def.cancel.signSum}
\lean{AlgebraicCombinatorics.SubtractiveMethods.signSum}
\leanhelper
\leanok
For a sign tuple $e : \operatorname{Fin} d \to \mathbb{Z}/2\mathbb{Z}$, the
\emph{sign sum} is the integer
\[
\operatorname{signSum}(e) := \sum_{k=1}^{d} \operatorname{toSign}(e_k).
\]
\end{definition}

\begin{definition}[Set of all-even tuples]
\label{def.cancel.allEvenTuples}
\lean{AlgebraicCombinatorics.SubtractiveMethods.allEvenTuples}
\leanhelper
\leanok
The set of all-even $n$-tuples in $[d]^n$ is
\[
\operatorname{allEvenTuples}(n, d)
:= \{x \in (\operatorname{Fin} n \to \operatorname{Fin} d) \mid x \text{ is all-even}\}.
\]
\end{definition}

\begin{theorem}
\label{thm.cancel.all-even}
\lean{AlgebraicCombinatorics.SubtractiveMethods.allEven_count_formula}
\leantarget
\leanok
Let $n\in\mathbb{N}$ and $d\in\mathbb{N}$.
Then, the number of all all-even $n$-tuples $\left(  x_{1},x_{2},\ldots
,x_{n}\right)  \in\left[  d\right]  ^{n}$ is
\[
\dfrac{1}{2^{d}}\sum_{k=0}^{d}\dbinom{d}{k}\left(  d-2k\right)  ^{n}.
\]
More precisely (avoiding division), the number of all-even $n$-tuples
multiplied by $2^d$ equals $\sum_{k=0}^{d}\binom{d}{k}(d-2k)^n$.
\end{theorem}

\begin{proof}
\leanok
Lemma~\ref{lem.cancel.all-even.l1} yields
\begin{align*}
&  \sum_{\left(  e_{1},e_{2},\ldots,e_{d}\right)  \in\left\{  1,-1\right\}
^{d}}\left(  e_{1}+e_{2}+\cdots+e_{d}\right)  ^{n}\\
&  =\sum_{\left(  x_{1},x_{2},\ldots,x_{n}\right)  \in\left[  d\right]  ^{n}%
}\ \ \sum_{\left(  e_{1},e_{2},\ldots,e_{d}\right)  \in\left\{  1,-1\right\}
^{d}}e_{x_{1}}e_{x_{2}}\cdots e_{x_{n}}.
\end{align*}
By Lemma~\ref{lem.cancel.sum-signProduct-eq-allEven-count}
(which splits according to whether each tuple is all-even or not,
using Lemma~\ref{lem.cancel.all-even.l2a} and
Lemma~\ref{lem.cancel.all-even.l2b}), the right-hand side equals
$(\text{\# of all-even tuples}) \cdot 2^d$.

On the other hand, a $d$-tuple $\left(  e_{1},e_{2},\ldots,e_{d}\right)  \in\left\{
1,-1\right\}  ^{d}$ is uniquely determined by the set $S = \left\{  i\in\left[
d\right]  \ \mid\ e_{i}=-1\right\}$, and for this tuple we have
$e_{1}+e_{2}+\cdots+e_{d}=d-2\left\vert S\right\vert$.
By Lemma~\ref{lem.cancel.sum-signSum-pow-eq-sum-subset}, we can rewrite
the sum over sign tuples as a sum over subsets.
By Lemma~\ref{lem.cancel.sum-subset-eq-sum-choose}, grouping by
cardinality $|S| = k$ yields
\[
\sum_{\left(  e_{1},e_{2},\ldots,e_{d}\right)  \in\left\{  1,-1\right\}  ^{d}%
}\left(  e_{1}+e_{2}+\cdots+e_{d}\right)  ^{n}=\sum_{k=0}^{d}\dbinom{d}%
{k}\left(  d-2k\right)  ^{n}.
\]
Comparing the two expressions gives
\[
\left(  \text{\# of all-even tuples}\right)  \cdot2^{d}=\sum_{k=0}^{d}\dbinom{d}%
{k}\left(  d-2k\right)  ^{n}.
\]
\end{proof}

\subsection{Lemmas for the proof}

\begin{lemma}
\label{lem.cancel.all-even.signProduct-eq-prod-pow}
\lean{AlgebraicCombinatorics.SubtractiveMethods.signProduct_eq_prod_pow}
\leanhelper
\leanok
Let $n,d\in\mathbb{N}$. Let $e = (e_1, \ldots, e_d) \in \{1,-1\}^d$ and
$x = (x_1, \ldots, x_n) \in [d]^n$. For each $k \in [d]$, let $m_k$ be
the multiplicity of $k$ in $x$. Then
\[
e_{x_{1}}e_{x_{2}}\cdots e_{x_{n}}=e_{1}^{m_{1}}e_{2}^{m_{2}}\cdots
e_{d}^{m_{d}}.
\]
\end{lemma}

\begin{proof}
\leanok
Each factor of the product $e_{x_{1}}e_{x_{2}}\cdots e_{x_{n}}$ is one of the $d$
entries $e_{1},e_{2},\ldots,e_{d}$. Moreover, each given
entry $e_{k}$ appears exactly $m_{k}$ times in the product $e_{x_{1}}e_{x_{2}%
}\cdots e_{x_{n}}$, because $k$ appears exactly $m_{k}$ times among the
subscripts $x_{1},x_{2},\ldots,x_{n}$. Thus the
product equals $e_{1}^{m_{1}}e_{2}^{m_{2}}\cdots e_{d}^{m_{d}}$.
\end{proof}

\begin{lemma}
\label{lem.cancel.all-even.l1}
\lean{AlgebraicCombinatorics.SubtractiveMethods.sum_signSum_pow_eq_sum_signProduct}
\leantarget
\leanok
Let $n,d\in\mathbb{N}$. Then,
\begin{align*}
&  \sum_{\left(  e_{1},e_{2},\ldots,e_{d}\right)  \in\left\{  1,-1\right\}
^{d}}\left(  e_{1}+e_{2}+\cdots+e_{d}\right)  ^{n}\\
&  =\sum_{\left(  x_{1},x_{2},\ldots,x_{n}\right)  \in\left[  d\right]  ^{n}%
}\ \ \sum_{\left(  e_{1},e_{2},\ldots,e_{d}\right)  \in\left\{  1,-1\right\}
^{d}}e_{x_{1}}e_{x_{2}}\cdots e_{x_{n}}.
\end{align*}
\end{lemma}

\begin{proof}
\leanok
For each $\left(  e_{1},e_{2},\ldots,e_{d}\right)  \in\left\{  1,-1\right\}
^{d}$, we have
\begin{align*}
\left(  e_{1}+e_{2}+\cdots+e_{d}\right)  ^{n}
&  =\left(  \sum_{k=1}^{d}e_{k}\right)  ^{n}
=\sum_{\left(  x_{1},x_{2},\ldots,x_{n}\right)  \in\left[
d\right]  ^{n}}e_{x_{1}}e_{x_{2}}\cdots e_{x_{n}}
\end{align*}
(by the generalized distributive law / product rule).
Summing over all $\left(
e_{1},e_{2},\ldots,e_{d}\right)  \in\left\{  1,-1\right\}  ^{d}$ and
interchanging the order of summation yields the result.
\end{proof}

\begin{lemma}
\label{lem.cancel.all-even.l2}
\lean{AlgebraicCombinatorics.SubtractiveMethods.sum_signProduct_not_allEven, AlgebraicCombinatorics.SubtractiveMethods.sum_signProduct_allEven}
\leantarget
\leanok
Let $n,d\in\mathbb{N}$. Let $\left(  x_{1},x_{2},\ldots,x_{n}\right)
\in\left[  d\right]  ^{n}$.

\textbf{(a)} If the $n$-tuple $\left(  x_{1},x_{2},\ldots,x_{n}\right)$ is
not all-even, then
$\sum_{\left(  e_{1},\ldots,e_{d}\right)  \in\left\{  1,-1\right\}  ^{d}}
e_{x_{1}}e_{x_{2}}\cdots e_{x_{n}}=0$.

\textbf{(b)} If the $n$-tuple $\left(  x_{1},x_{2},\ldots,x_{n}\right)$ is
all-even, then
$\sum_{\left(  e_{1},\ldots,e_{d}\right)  \in\left\{  1,-1\right\}  ^{d}}
e_{x_{1}}e_{x_{2}}\cdots e_{x_{n}}=2^{d}$.
\end{lemma}

\begin{proof}
\leanok
This combines parts \textbf{(a)} and \textbf{(b)}.
\end{proof}

\subsection{Involution argument for part (a)}

\begin{definition}
\label{def.cancel.flipSign}
\lean{AlgebraicCombinatorics.SubtractiveMethods.flipSign}
\leanhelper
\leanok
For $k \in [d]$, the map $\operatorname{flipSign}_k : \{1,-1\}^d \to \{1,-1\}^d$
sends a $d$-tuple $(e_1, \ldots, e_d)$ to the tuple obtained by flipping
the sign of the $k$-th entry: i.e., replacing $e_k$ by $-e_k$ and
leaving all other entries unchanged.
\end{definition}

\begin{lemma}
\label{lem.cancel.flipSign-involutive}
\lean{AlgebraicCombinatorics.SubtractiveMethods.flipSign_involutive}
\leanhelper
\leanok
For any $k \in [d]$, the map $\operatorname{flipSign}_k$ is an involution.
\end{lemma}

\begin{proof}
\leanok
Flipping the sign of the $k$-th entry twice restores the original sign.
\end{proof}

\begin{lemma}
\label{lem.cancel.flipSign-ne}
\lean{AlgebraicCombinatorics.SubtractiveMethods.flipSign_ne}
\leanhelper
\leanok
For any $k \in [d]$ and $e \in \{1,-1\}^d$, we have
$\operatorname{flipSign}_k(e) \neq e$.
\end{lemma}

\begin{proof}
\leanok
Since $e_k \in \{1, -1\}$, flipping the sign of $e_k$ produces a different
value ($-e_k \neq e_k$), so the resulting tuple differs from the original.
\end{proof}

\begin{lemma}[flipSign at the flipped position]
\label{lem.cancel.flipSign-self}
\lean{AlgebraicCombinatorics.SubtractiveMethods.flipSign_self}
\leanhelper
\leanok
For any $k \in [d]$ and $e \in \{1,-1\}^d$,
$(\operatorname{flipSign}_k(e))_k = e_k + 1$ (in $\mathbb{Z}/2\mathbb{Z}$).
\end{lemma}

\begin{proof}
\leanok
By definition of $\operatorname{flipSign}$, which updates position $k$.
\end{proof}

\begin{lemma}[flipSign leaves other positions unchanged]
\label{lem.cancel.flipSign-ne-self}
\lean{AlgebraicCombinatorics.SubtractiveMethods.flipSign_ne_self}
\leanhelper
\leanok
For any $k, j \in [d]$ with $j \neq k$ and $e \in \{1,-1\}^d$,
$(\operatorname{flipSign}_k(e))_j = e_j$.
\end{lemma}

\begin{proof}
\leanok
The update only affects position $k$, so all other positions are unchanged.
\end{proof}

\begin{lemma}
\label{lem.cancel.signProduct-flipSign}
\lean{AlgebraicCombinatorics.SubtractiveMethods.signProduct_flipSign}
\leanhelper
\leanok
Let $e \in \{1,-1\}^d$, $x \in [d]^n$, and $k \in [d]$ such that the
multiplicity $m_k$ of $k$ in $x$ is odd. Then
\[
\operatorname{signProduct}(\operatorname{flipSign}_k(e), x) =
-\operatorname{signProduct}(e, x).
\]
\end{lemma}

\begin{proof}
\leanok
Using the product-of-powers representation (Lemma~\ref{lem.cancel.all-even.signProduct-eq-prod-pow}),
flipping the $k$-th entry multiplies the product by $(-1)^{m_k} = -1$
(since $m_k$ is odd), while all other factors remain unchanged.
\end{proof}

\begin{lemma}
\label{lem.cancel.all-even.l2a}
\lean{AlgebraicCombinatorics.SubtractiveMethods.sum_signProduct_not_allEven}
\leanhelper
\leanok
Let $n,d\in\mathbb{N}$. Let $\left(  x_{1}%
,x_{2},\ldots,x_{n}\right)  \in\left[  d\right]  ^{n}$.
If the $n$-tuple $\left(  x_{1},x_{2},\ldots,x_{n}\right)  $ is
not all-even, then
\[
\sum_{\left(  e_{1},e_{2},\ldots,e_{d}\right)  \in\left\{  1,-1\right\}  ^{d}%
}e_{x_{1}}e_{x_{2}}\cdots e_{x_{n}}=0.
\]
\end{lemma}

\begin{proof}
\leanok
Since the tuple is not all-even, there exists some
$k\in\left[  d\right]  $ such that $m_{k}$ is odd.

Define $f:\left\{  1,-1\right\}  ^{d}\rightarrow\left\{  1,-1\right\}  ^{d}$
as $\operatorname{flipSign}_k$ (flipping the sign of the $k$-th entry).
By Lemma~\ref{lem.cancel.flipSign-involutive}, $f$ is an involution.
By Lemma~\ref{lem.cancel.flipSign-ne}, $f$ has no fixed points.
By Lemma~\ref{lem.cancel.signProduct-flipSign},
$\operatorname{signProduct}(f(e), x) = -\operatorname{signProduct}(e, x)$.

Hence, the terms cancel in pairs, giving
\[
\sum_{\left(  e_{1},e_{2},\ldots,e_{d}\right)  \in\left\{  1,-1\right\}  ^{d}%
}e_{x_{1}}e_{x_{2}}\cdots e_{x_{n}}=0.
\]
The Lean proof uses the involution summation lemma from Mathlib.
\end{proof}

\begin{lemma}
\label{lem.cancel.all-even.l2b}
\lean{AlgebraicCombinatorics.SubtractiveMethods.sum_signProduct_allEven}
\leanhelper
\leanok
Let $n,d\in\mathbb{N}$. Let $\left(  x_{1}%
,x_{2},\ldots,x_{n}\right)  \in\left[  d\right]  ^{n}$.
If the $n$-tuple $\left(  x_{1},x_{2},\ldots,x_{n}\right)  $ is
all-even, then
\[
\sum_{\left(  e_{1},e_{2},\ldots,e_{d}\right)  \in\left\{  1,-1\right\}  ^{d}%
}e_{x_{1}}e_{x_{2}}\cdots e_{x_{n}}=2^{d}.
\]
\end{lemma}

\begin{proof}
\leanok
Since the tuple is all-even, all the multiplicities $m_{1},m_{2},\ldots,m_{d}$
are even.

Let $\left(  e_{1},e_{2},\ldots,e_{d}\right)  \in\left\{  1,-1\right\}
^{d}$. For each $k\in\left[  d\right]  $, we have $e_{k}\in\left\{  1,-1\right\}  $
and $m_{k}$ is even, so $e_{k}^{m_{k}}=1$.
By Lemma~\ref{lem.cancel.all-even.signProduct-eq-prod-pow},
$e_{x_{1}}e_{x_{2}}\cdots e_{x_{n}} = e_{1}^{m_{1}}e_{2}^{m_{2}}\cdots e_{d}^{m_{d}} = 1$.

Since every sign product equals $1$ and there are $2^d$ sign tuples,
\[
\sum_{\left(  e_{1},e_{2},\ldots,e_{d}\right)  \in\left\{  1,-1\right\}  ^{d}%
}e_{x_{1}}e_{x_{2}}\cdots e_{x_{n}}=2^{d}.
\]
\end{proof}

\begin{lemma}
\label{lem.cancel.sum-signProduct-eq-allEven-count}
\lean{AlgebraicCombinatorics.SubtractiveMethods.sum_signProduct_eq_allEven_count}
\leanhelper
\leanok
Let $n,d\in\mathbb{N}$. Then
\[
\sum_{\left(  x_{1},x_{2},\ldots,x_{n}\right)  \in\left[  d\right]  ^{n}%
}\ \ \sum_{\left(  e_{1},e_{2},\ldots,e_{d}\right)  \in\left\{  1,-1\right\}
^{d}}e_{x_{1}}e_{x_{2}}\cdots e_{x_{n}}
= \left(\text{\# of all-even tuples in } [d]^n\right) \cdot 2^d.
\]
\end{lemma}

\begin{proof}
\leanok
Split the outer sum according to whether each tuple is all-even or not.
By Lemma~\ref{lem.cancel.all-even.l2a}, the non-all-even tuples contribute $0$.
By Lemma~\ref{lem.cancel.all-even.l2b}, each all-even tuple contributes $2^d$.
Hence the total is $(\text{\# of all-even tuples}) \cdot 2^d$.
\end{proof}

\subsection{Computing the left-hand side}

\begin{lemma}
\label{lem.cancel.signSum-eq-card}
\lean{AlgebraicCombinatorics.SubtractiveMethods.signSum_eq_card}
\leanhelper
\leanok
For a sign tuple $e \in \{1,-1\}^d$, let $S = \{i \in [d] \mid e_i = -1\}$.
Then
\[
e_1 + e_2 + \cdots + e_d = d - 2|S|.
\]
\end{lemma}

\begin{proof}
\leanok
Each entry $e_i$ equals $-1$ if $i \in S$ and $1$ if $i \notin S$.
The sum has $|S|$ addends equal to $-1$ and $d - |S|$ addends equal to $1$,
giving $|S| \cdot (-1) + (d - |S|) \cdot 1 = d - 2|S|$.
\end{proof}

\begin{definition}[Sign tuple to subset]
\label{def.cancel.signTupleToSubset}
\lean{AlgebraicCombinatorics.SubtractiveMethods.signTupleToSubset}
\leanhelper
\leanok
The map $\operatorname{signTupleToSubset} : (\operatorname{Fin} d \to \mathbb{Z}/2\mathbb{Z}) \to \mathcal{P}(\operatorname{Fin} d)$
sends a sign tuple $e$ to the set $S = \{k \in \operatorname{Fin} d \mid e_k = 1\}$
(i.e., the positions where $\operatorname{toSign}(e_k) = -1$).
\end{definition}

\begin{definition}[Subset to sign tuple]
\label{def.cancel.subsetToSignTuple}
\lean{AlgebraicCombinatorics.SubtractiveMethods.subsetToSignTuple}
\leanhelper
\leanok
The inverse map $\operatorname{subsetToSignTuple} : \mathcal{P}(\operatorname{Fin} d) \to (\operatorname{Fin} d \to \mathbb{Z}/2\mathbb{Z})$
sends a subset $S \subseteq \operatorname{Fin} d$ to the sign tuple where
position $k$ has value $1$ (representing $-1$) if $k \in S$, and $0$
(representing $+1$) otherwise.
\end{definition}

\begin{lemma}
\label{lem.cancel.signTupleEquivSubset}
\lean{AlgebraicCombinatorics.SubtractiveMethods.signTupleEquivSubset}
\leanhelper
\leanok
There is a bijection between sign tuples $\{1,-1\}^d$ and subsets of $[d]$,
sending each tuple $e$ to the set $S = \{i \in [d] \mid e_i = -1\}$.
\end{lemma}

\begin{proof}
\leanok
The inverse sends a subset $S \subseteq [d]$ to the sign tuple where
position $k$ has $e_k = -1$ if $k \in S$ and $e_k = 1$ otherwise.
\end{proof}

\begin{lemma}[Sign sum under the subset correspondence]
\label{lem.cancel.signSum-subsetToSignTuple}
\lean{AlgebraicCombinatorics.SubtractiveMethods.signSum_subsetToSignTuple}
\leanhelper
\leanok
For a subset $S \subseteq [d]$, the sign sum of the corresponding sign tuple is
\[
\operatorname{signSum}(\operatorname{subsetToSignTuple}(S)) = d - 2|S|.
\]
\end{lemma}

\begin{proof}
\leanok
By Lemma~\ref{lem.cancel.signSum-eq-card}, the sign sum equals $d - 2|T|$
where $T = \operatorname{signTupleToSubset}(\operatorname{subsetToSignTuple}(S)) = S$.
\end{proof}

\begin{lemma}
\label{lem.cancel.sum-signSum-pow-eq-sum-subset}
\lean{AlgebraicCombinatorics.SubtractiveMethods.sum_signSum_pow_eq_sum_subset}
\leanhelper
\leanok
Let $n,d\in\mathbb{N}$. Then
\[
\sum_{\left(  e_{1},e_{2},\ldots,e_{d}\right)  \in\left\{  1,-1\right\}  ^{d}%
}\left(  e_{1}+e_{2}+\cdots+e_{d}\right)  ^{n}=\sum_{S\subseteq\left[
d\right]  }\left(  d-2\left\vert S\right\vert \right)  ^{n}.
\]
\end{lemma}

\begin{proof}
\leanok
Substitute via the bijection of Lemma~\ref{lem.cancel.signTupleEquivSubset}
and use Lemma~\ref{lem.cancel.signSum-eq-card}.
\end{proof}

\begin{lemma}
\label{lem.cancel.sum-subset-eq-sum-choose}
\lean{AlgebraicCombinatorics.SubtractiveMethods.sum_subset_eq_sum_choose}
\leanhelper
\leanok
Let $n,d\in\mathbb{N}$. Then
\[
\sum_{S\subseteq\left[  d\right]  }\left(  d-2\left\vert S\right\vert
\right)  ^{n} = \sum_{k=0}^{d}\dbinom{d}{k}\left(  d-2k\right)  ^{n}.
\]
\end{lemma}

\begin{proof}
\leanok
Split the sum over subsets $S$ according to the value of $|S| = k$.
For each $k$, all subsets of size $k$ contribute the same summand $(d-2k)^n$,
and the number of $k$-element subsets of $[d]$ is $\binom{d}{k}$.
\end{proof}

\subsection{Helper lemmas for toSign}

\begin{lemma}
\label{lem.cancel.toSign-sq}
\lean{AlgebraicCombinatorics.SubtractiveMethods.toSign_sq}
\leanhelper
\leanok
For any $b \in \mathbb{Z}/2\mathbb{Z}$, we have $(\operatorname{toSign} b)^2 = 1$.
\end{lemma}

\begin{proof}
\leanok
By case analysis: both $1^2 = 1$ and $(-1)^2 = 1$.
\end{proof}

\begin{lemma}
\label{lem.cancel.toSign-pow-even}
\lean{AlgebraicCombinatorics.SubtractiveMethods.toSign_pow_even}
\leanhelper
\leanok
For any $b \in \mathbb{Z}/2\mathbb{Z}$ and even $m \in \mathbb{N}$,
we have $(\operatorname{toSign} b)^m = 1$.
\end{lemma}

\begin{proof}
\leanok
Write $m = 2k$. Then $(\operatorname{toSign} b)^m = ((\operatorname{toSign} b)^2)^k = 1^k = 1$.
\end{proof}

\begin{lemma}[Odd powers of a sign]
\label{lem.cancel.toSign-pow-odd}
\lean{AlgebraicCombinatorics.SubtractiveMethods.toSign_pow_odd}
\leanhelper
\leanok
For any $b \in \mathbb{Z}/2\mathbb{Z}$ and odd $m \in \mathbb{N}$,
we have $(\operatorname{toSign} b)^m = \operatorname{toSign} b$.
\end{lemma}

\begin{proof}
\leanok
Write $m = 2k + 1$. Then
$(\operatorname{toSign} b)^m = ((\operatorname{toSign} b)^2)^k \cdot \operatorname{toSign} b
= 1^k \cdot \operatorname{toSign} b = \operatorname{toSign} b$.
\end{proof}

\begin{lemma}
\label{lem.cancel.toSign-add-one}
\lean{AlgebraicCombinatorics.SubtractiveMethods.toSign_add_one}
\leanhelper
\leanok
For any $b \in \mathbb{Z}/2\mathbb{Z}$, we have
$\operatorname{toSign}(b+1) = -\operatorname{toSign}(b)$.
\end{lemma}

\begin{proof}
\leanok
By case analysis on $b \in \{0, 1\}$: adding $1$ in $\mathbb{Z}/2\mathbb{Z}$ toggles between $0$ and $1$, and $\operatorname{toSign}(0) = 1 = -(-1) = -\operatorname{toSign}(1)$ while $\operatorname{toSign}(1) = -1 = -(1) = -\operatorname{toSign}(0)$.
\end{proof}

\subsection{Corollaries of the counting formula}

\begin{corollary}[Divisibility by $2^d$]
\label{cor.cancel.sum-choose-pow-dvd}
\lean{AlgebraicCombinatorics.SubtractiveMethods.sum_choose_pow_dvd_pow_two}
\leanhelper
\leanok
For any $n, d \in \mathbb{N}$, we have
\[
2^d \mid \sum_{k=0}^{d} \binom{d}{k} (d - 2k)^n.
\]
\end{corollary}

\begin{proof}
\leanok
By Theorem~\ref{thm.cancel.all-even}, the sum equals
$(\text{\# of all-even tuples}) \cdot 2^d$, which is divisible by $2^d$.
\end{proof}

\begin{corollary}[Nonnegativity]
\label{cor.cancel.sum-choose-pow-nonneg}
\lean{AlgebraicCombinatorics.SubtractiveMethods.sum_choose_pow_nonneg}
\leanhelper
\leanok
For any $n, d \in \mathbb{N}$, we have
\[
\sum_{k=0}^{d} \binom{d}{k} (d - 2k)^n \geq 0.
\]
This is not obvious from the formula when $n$ is odd, since the summands
can be negative.
\end{corollary}

\begin{proof}
\leanok
By Theorem~\ref{thm.cancel.all-even}, the sum equals
$(\text{\# of all-even tuples}) \cdot 2^d$, which is a product of two nonnegative
integers.
\end{proof}
